\section{Lifting on the divisor}\label{sec:integral}

The divisor $\Spec S\subset\Spec A$ now supplies an integral class
with zero generic restriction. We first lift a coefficient class
to $K_3(D)$ and then extend that lift across the closed point of $S$.
Its coefficient image will be tested for nonvanishing only after
this construction.

Let $M_5$ be the mod-five Moore spectrum. We use
$K_i(B;\Z/5)=\pi_i(K(B)\wedge M_5)$, and the same convention for
$G$-theory. The coefficient sequence is
\begin{equation}\label{eq:coeff}
 K_i(B)\xrightarrow{5}K_i(B)\xrightarrow{\rho_B}K_i(B;\Z/5)
 \xrightarrow{\partial_{\mathrm{coef}}}K_{i-1}(B)
 \xrightarrow{5}K_{i-1}(B).
\end{equation}
Here $\rho_B$ denotes coefficient reduction, and
$\partial_{\mathrm{coef}}$ denotes the connecting map.
Multiplication of a coefficient class by an integral class uses
the natural $K(B)$-module structure on $K(B)\wedge M_5$.

For an element $u$ of a field's multiplicative group, write $[u]$
for its integral $K_1$ class. Since
$[\zeta]\in K_1(L)=L^\times$ has order five, exactness permits
us to fix
\[
 \beta\in K_2(L;\Z/5),\qquad
 \partial_{\mathrm{coef}}\beta=[\zeta].
\]
For every finite extension $E/L$, write $\beta_E$ for its restriction.
The class to be lifted is
\[
 a_D=\beta_D\cdot[t]\in K_3(D;\Z/5).
\]

Resolution identifies $K$-theory with coherent
$G$-theory for the regular rings $S$ and $A$; closed-immersion
pushforward is taken in $G$-theory
\cite[\S7.1 and \S7, paragraph 2.7 and the discussion following (2.8)]{QuillenHigher}.

\begin{proposition}\label{prop:integral}
There exists $b\in K_3(S)$ whose restriction to $D$, followed
by coefficient reduction, is $a_D$. If
$j:\Spec S\hookrightarrow\Spec A$ and $c=j_*b\in G_3(A)=K_3(A)$,
then $c|_{\Frac A}=0$. Moreover, for
$T=\Spec A[1/\pi]$ and $i:\Spec D\hookrightarrow T$,
\[
 c|_T\bmod5=i_*a_D\quad\text{in }G_3(T;\Z/5).
\]
\end{proposition}

\begin{proof}
The coefficient cofiber sequence is a sequence of $K(D)$-modules.
Its connecting map is therefore compatible with multiplication by
an integral $K(D)$ class, up to the suspension sign. By \eqref{eq:norm-t},
$\zeta=\eta^5$ and $t=1-\eta$. Both $\eta$ and $1-\eta$
are nonzero: the first has fifth power $\zeta\ne0$, and
the second is the uniformizer $t$ of $S$. Matsumoto's
Steinberg relation \cite[\S5, Lemma 5.6, Theorem 5.10, and
Corollary 5.11]{Matsumoto}, with the identification of symbols and
products in Quillen $K_2(D)$
\cite[III.6.1, IV.1.7.1, IV.1.10.1, and IV.7.2]{Weibel}, gives
\begin{equation}\label{eq:steinberg}
 \partial_{\mathrm{coef}}a_D
 =\pm[\zeta]\cdot[t]
 =\pm5\{\eta,1-\eta\}=0.
\end{equation}
The sign does not affect the vanishing. Exactness of
\eqref{eq:coeff} yields an integral lift
$\widetilde a_D\in K_3(D)$ of $a_D$.
The boundary $[\zeta]\in K_1(D)[5]$ of $\beta_D$ itself is
still nonzero, even though $\zeta=\eta^5$. Being a fifth power
makes the class of $\zeta$ vanish in $D^\times/D^{\times5}$,
not in the integral group $K_1(D)=D^\times$.
The Steinberg relation is what kills the boundary of the product
$a_D$.

Localization and d\'evissage for the DVR $S$ give
\cite[\S7, Proposition 3.2]{QuillenHigher}
\[
 K_3(S)\longrightarrow K_3(D)\longrightarrow K_2(\F_5).
\]
The last group is zero by Quillen's finite-field computation
\cite{QuillenFinite}. We may therefore choose
$b\in K_3(S)$ restricting to $\widetilde a_D$. This is a
second lift, distinct from the coefficient lift to $K_3(D)$.

Let $c=j_*b\in G_3(A)=K_3(A)$. Its restriction to
$\Spec A[1/y]$ is zero because the closed immersion $j$
has empty pullback there. The element $y$ is nonzero in the
domain $A$: its quotient $S$ has dimension one, whereas $A$
has dimension two. Hence $y$ becomes invertible in $\Frac A$,
and $c|_{\Frac A}=0$ integrally.

Base change of $j$ along $T\to\Spec A$ is
$i:\Spec D\hookrightarrow T$, since $S[1/\pi]=D$.
Flat base change for coherent pushforward
\cite[V.3.7.2]{Weibel} here is the natural isomorphism
\[
 A[1/\pi]\otimes_A M\cong D\otimes_S M
\]
for a finite $S$-module $M$, with both sides viewed as
$A[1/\pi]$-modules. Thus restriction of scalars commutes with
localization. Naturality of coefficient reduction gives the
commutative diagram
\[
\begin{array}{ccc}
 K_3(S)&\xrightarrow{\ j_*\ }&G_3(A)\\
 \big\downarrow&&\big\downarrow\\
 K_3(D;\Z/5)&\xrightarrow{\ i_*\ }&G_3(T;\Z/5).
\end{array}
\]
Both vertical arrows mean restriction followed by coefficient
reduction. The left arrow sends $b$ to $a_D$,
which proves the stated identity.
\end{proof}
